\section{离散傅里叶变换}\label{sec:DFT}

\subsection{离散傅里叶变换的引入：离散傅里叶级数}\label{sub:intro_DFT}

沿着\refsub{dtft}的思路，我们仿照傅里叶级数来进行周期序列的频谱分析。设$x[n]$是周期为$N$的序列，即$x[n+N]=x[n],\ \forall n\in\mathbb{Z}$，考虑将其表示为一组复指数序列的线性组合，则每个复指数序列也应该以$N$为周期，即
\begin{equation}
    E_{\Omega}[n]=\e^{\im\Omega n},\quad E_{\Omega}[n+N]=E_{\Omega}[n],\quad \forall n\in\mathbb{Z}.
\end{equation}
\refsub{digital_freq}中已经指出，复指数序列的数字角频率必须取值于有限的集合：
\begin{equation}
    \braces{\Omega_k=\frac{2\pi k}{N}\;\middle|\;k=0,1,\dots,N-1}.
\end{equation}
这个集合恰有$N$个元素。简便起见，这里数字角频率的范围取为$[0,2\pi)$。换言之，我们希望找到一个与$x[n]$具有相同周期的序列$X[k]$，使得
\begin{equation}
    x[n]=\frac{1}{N}\sum_{k=0}^{N-1}X[k]E_{\Omega_k}[n]=\frac{1}{N}\sum_{k=0}^{N-1}X[k]\e^{\im\frac{2\pi}{N}kn}
\end{equation}
（我们马上会看到系数$1/N$的作用）。

为了验证这个想法的可行性，将以上公式写为矩阵形式：
\begin{equation}
    \begin{bmatrix}
        x[0]   \\
        x[1]   \\
        \vdots \\
        x[N-1]
    \end{bmatrix}=\frac{1}{N}
    \begin{bmatrix}
        1      & 1                                  & \cdots & 1                                       \\
        1      & \e^{\im\frac{2\pi}{N}}             & \cdots & \e^{\im\frac{2\pi}{N}(N-1)}             \\
        \vdots & \vdots                             & \ddots & \vdots                                  \\
        1      & \e^{\im\frac{2\pi}{N}(N-1)}        & \cdots & \e^{\im\frac{2\pi}{N}(N-1)^2}
    \end{bmatrix}
    \begin{bmatrix}
        X[0]   \\
        X[1]   \\
        \vdots \\
        X[N-1]
    \end{bmatrix}.
\end{equation}

习惯上，将$\e^{-\im\frac{2\pi}{N}}$记为$W_N$，则上式中的矩阵可以写为
\begin{equation}
    \mathcal{F}=\begin{bmatrix}
        1      & 1            & \cdots & 1            \\
        1      & W_N^{-1}     & \cdots & W_N^{-(N-1)} \\
        \vdots & \vdots       & \ddots & \vdots       \\
        1      & W_N^{-(N-1)} & \cdots & W_N^{-(N-1)^2}
    \end{bmatrix}.
\end{equation}

这是一个\textbf{范德蒙德矩阵}(Vandermonde Matrix)，它的行列式为
\begin{equation}
    \det\mathcal{F}=\prod_{0\leq j<k\leq N-1}(W_N^{-k}-W_N^{-j})\neq 0.
\end{equation}

矩阵$\mathcal{F}$可逆，因此$X[k]$是良定义的，即存在唯一的$X[k]$使得
\begin{equation}
    x[n]=\frac{1}{N}\sum_{k=0}^{N-1}X[k]W_N^{-kn},
\end{equation}
并且可以通过逆矩阵求得$X[k]$的表达式。可以验证，$\mathcal{F}$为正交矩阵：
\begin{equation}
    \mathcal{F}\mathcal{F}^H=N\cdot I,
\end{equation}
其中$\mathcal{F}^H$为$\mathcal{F}$的共轭转置矩阵，因此
\begin{equation}
    \mathcal{F}^{-1}=\frac{1}{N}\mathcal{F}^H=\frac{1}{N}
    \begin{bmatrix}
        1      & 1         & \cdots & 1           \\
        1      & W_N       & \cdots & W_N^{N-1}   \\
        \vdots & \vdots    & \ddots & \vdots      \\
        1      & W_N^{N-1} & \cdots & W_N^{(N-1)^2}
    \end{bmatrix}.
\end{equation}

这样，我们就得到了\textbf{离散傅里叶级数}(Discrete Fourier Series, DFS)的定义：
\begin{definition}[离散傅里叶级数]
    设$x[n]$是周期为$N$的序列，则$x[n]$的DFS定义为
    \begin{equation}
        x[n]=\frac{1}{N}\sum_{k=0}^{N-1}X[k]W_N^{-kn},
    \end{equation}
    其中$W_N=\e^{\im\frac{2\pi}{N}}$，$X[k]=\sum_{n=0}^{N-1}x[n]W_N^{-kn}$。
\end{definition}

尽管在分析序列时没有“高次谐波分量”，但从以上讨论可以看出，只要有这$N$个离散频率成分，就可以完全表征周期为$N$的序列的频谱特性。至此，我们完美地回答了“如何对序列进行分解与合成”的问题。

本小节的主要目标——\textbf{离散傅里叶变换}(Discrete Fourier Transform, DFT)就是求解DFS的系数的这个变换，只是习惯上认为离散傅里叶变换是针对于\textbf{有限长序列}（或者说仅在有限长度的区间上非零的序列）而言的\cite{dsp}。\textbf{本书不区分有限长序列与周期序列}，出于以下两个考虑：其一，它们之间的互相转化非常简单；其二，这样的写法更有利于展示离散傅里叶变换的一系列运算性质的本质。

上述内容归结为以下的定义和定理：

\begin{definition}[离散傅里叶变换]\label{def:DFT}
    设$x[n]$是长为$N$的序列，则$x[n]$的DFT定义为
    \begin{equation}
        X[k]=\sum_{n=0}^{N-1}x[n]W_N^{kn},\quad k=0,1,\dots,N-1,
    \end{equation}
    记$x[n]\xleftrightarrow{\mathrm{DFT}}X[k]$，$X[k]=\DFT(x[n])$。
\end{definition}

\begin{theorem}[离散傅里叶反演公式]\label{thm:DFT_inv}
    设$x[n]$是长为$N$的序列，其离散傅里叶变换为$X[k]$，则$x[n]$可以由$X[k]$通过离散傅里叶逆变换(IDFT)还原：
    \begin{equation}
        x[n]=\DFT^{-1}(X[k])=\frac{1}{N}\sum_{k=0}^{N-1}X[k]W_N^{-kn},\quad n=0,1,\dots,N-1.
    \end{equation}
\end{theorem}

\begin{example}
    设$\mathbf{x}=[1,2,3,4]^T$，求其DFT。

    解：由定义，
    \begin{align}
        \begin{bmatrix}
            X[0] \\
            X[1] \\
            X[2] \\
            X[3]
        \end{bmatrix}&=\begin{bmatrix}
            1 & 1     & 1     & 1     \\
            1 & W_4   & W_4^2 & W_4^3 \\
            1 & W_4^2 & W_4^4 & W_4^6 \\
            1 & W_4^3 & W_4^6 & W_4^9
        \end{bmatrix}\begin{bmatrix}
            1 \\
            2 \\
            3 \\
            4
        \end{bmatrix} \nonumber \\
        &=\begin{bmatrix}
            1 & 1    & 1  & 1    \\
            1 & -\im & -1 & \im  \\
            1 & -1   & 1  & -1   \\
            1 & \im  & -1 & -\im
        \end{bmatrix}\begin{bmatrix}
            1 \\
            2 \\
            3 \\
            4
        \end{bmatrix}=\begin{bmatrix}
            10      \\
            -2+2\im \\
            -2      \\
            -2-2\im
        \end{bmatrix}.
    \end{align}
\end{example}

对于周期序列，其周期起点任意选取，都不会影响DFT及IDFT的结果。
\begin{definition}[$(\mathcal{P},\mathcal{Q})$-DFT]
    设$x[n]$是以$N$为周期的序列，$\mathcal{P},\mathcal{Q}$为任意的$N$个连续整数构成的指标集，即
    \begin{equation}
        \mathcal{P}=\{p,p+1,\cdots,p+N-1\},\quad \mathcal{Q}=\{q,q+1,\cdots,q+N-1\},
    \end{equation}
    则$x[n]$的$(\mathcal{P},\mathcal{Q})$-DFT定义为
    \begin{equation}
        X_{\mathcal{P},\mathcal{Q}}[k]=\sum_{n\in\mathcal{P}}x[n]W_N^{kn},\quad k\in\mathcal{Q}.
    \end{equation}
\end{definition}
\begin{proposition}
    $(\mathcal{P},\mathcal{Q})$-DFT的定义不依赖于$\mathcal{P},\mathcal{Q}$的选取，即
    \begin{equation}
        X_{\mathcal{P},\mathcal{Q}}[k]=X[k],\quad \forall \mathcal{P},\mathcal{Q}.
    \end{equation}
\end{proposition}
\begin{proof}
    记$p=rN+s$，则$\forall k\in\mathbb{Z}$，有
    \begin{align}
        X_{\mathcal{P},\mathcal{Q}}[k]&=\sum_{n\in\mathcal{P}}x[n]W_N^{kn} \nonumber \\
        &=\sum_{m=p}^{p+N-1}x[m]W_N^{km} \nonumber \\
        &=\sum_{m=0}^{N-1}x[m+s]W_N^{k(m+s)} \nonumber \\
        &=\sum_{m=s}^{N-1}x[m]W_N^{km}+\sum_{m=N}^{N+s-1}x[m]W_N^{km} \nonumber \\
        &=\sum_{m=s}^{N-1}x[m]W_N^{km}+\sum_{m=0}^{s-1}x[m]W_N^{km} \nonumber \\
        &=\sum_{m=0}^{N-1}x[m]W_N^{km} \nonumber \\
        &=X[k].
    \end{align}
    限制$k\in\mathcal{Q}$即得$X_{\mathcal{P},\mathcal{Q}}[k]=X[k]$。
\end{proof}

在此意义下，周期序列的向量表示不唯一，我们约定：周期序列对应的向量总是以$n=0$时刻的值作为第一个元素。

\subsection{离散傅里叶变换的运算性质}\label{sub:property_of_DFT}

\begin{proposition}[离散傅里叶变换的运算性质]
    \label{prop:DFT_operation}
    设$x[n],y[n]$是周期为$N$的序列，其DFT分别为$X[k],Y[k]$，则有：
    \begin{itemize}
        \item 线性性：$\forall a,b\in\mathbb{C}$，有
        \begin{equation}
            a x[n]+b y[n]\xleftrightarrow{\mathrm{DFT}} a X[k]+b Y[k]
        \end{equation}
        \item 周期性：$X[k+N]=X[k],Y[k+N]=Y[k]$为周期为$N$的序列
        \item 时移性质：$x[n-n_0]\xleftrightarrow{\mathrm{DFT}}\e^{-\im\frac{2\pi}{N}kn_0}X[k]=W_N^{kn_0}X[k]$
        \item 频移性质：$\e^{\im\frac{2\pi}{N}k_0 n}x[n]=W_N^{-k_0 n}x[n]\xleftrightarrow{\mathrm{DFT}}X[k-k_0]$
        \item 共轭性质：$\overline{x[n]}\xleftrightarrow{\mathrm{DFT}}\overline{X[-k]}$
    \end{itemize}
\end{proposition}

以上性质全都可以直接带入定义验证，这里不再赘述。

我们曾定义过序列的卷积和。类似\cref{def:period_conv}，对于周期序列，卷积和的定义也必须修改，只是此时一般不会添加$1/N$的系数：
\begin{definition}[周期序列的卷积和]
    设$x[n],y[n]$是周期为$N$的序列，则它们的卷积和定义为
    \begin{equation}
        (x*y)[n]=\sum_{m=0}^{N-1}x[m]y[n-m].
    \end{equation}
\end{definition}
为了区分两种卷积和，我们将非周期序列的卷积和称为\textbf{线性卷积}，而将周期序列的卷积和称为\textbf{循环卷积}或\textbf{圆卷积}。线性卷积的命名是由于它对参与运算的两序列都是线性的，而循环卷积的线性性必须保证序列具有相同的周期。一些文献中，也将循环卷积记作$\circledast$，或者用$\circledast_N$表示周期为$N$的序列的循环卷积。

根据离散傅里叶变换的周期性，时域和频域都是周期的，下面就来表述和验证卷积和定理：
\begin{theorem}[离散傅里叶变换的卷积和定理]\label{thm:DFT_conv}
    设$x[n],y[n]$是周期为$N$的序列，其DFT分别为$X[k],Y[k]$，则有
    \begin{align}
        (x*y)[n]&\xleftrightarrow{\mathrm{DFT}}X[k]Y[k], \\
        x[n]y[n]&\xleftrightarrow{\mathrm{DFT}}\frac{1}{N}(X * Y)[k].
    \end{align}
\end{theorem}
\begin{proof}
    \begin{align}
        X[k]Y[k]&=\paren{\sum_{n=0}^{N-1}x[n]W_N^{kn}}\paren{\sum_{m=0}^{N-1}y[m]W_N^{km}} \nonumber \\
        &=\sum_{n=0}^{N-1}\sum_{m=0}^{N-1}x[n]y[m]W_N^{k(n+m)} \nonumber \\
        &=\sum_{l=0}^{N-1}\paren{\sum_{n=0}^{N-1}x[n]y[l-n]}W_N^{kl} \qquad (l=n+m) \nonumber \\
        &=\DFT(x*y)[k].
    \end{align}

    可以仿照\refsub{convolution}，使用时域卷积定理给出频域卷积定理的证明；也可以直接带入定义验证：
    \begin{align}
        \frac{1}{N}(X * Y)[k]&=\frac{1}{N}\sum_{l=0}^{N-1}X[l]Y[k-l] \nonumber \\
        &=\frac{1}{N}\sum_{l=0}^{N-1}\paren{\sum_{n=0}^{N-1}x[n]W_N^{ln}}\paren{\sum_{m=0}^{N-1}y[m]W_N^{k m}W_N^{-lm}} \nonumber \\
        &=\frac{1}{N}\sum_{n=0}^{N-1}\sum_{m=0}^{N-1}x[n]y[m]W_N^{k m}\paren{\sum_{l=0}^{N-1}W_N^{l(n-m)}} \nonumber \\
        &=\sum_{n=0}^{N-1}x[n]y[n]W_N^{k n}.
    \end{align}
    最后一个等号用到了\refsub{z_operation}中引入的离散选择函数：
    \begin{equation}
        \shah_N[m-n]=\frac{1}{N}\sum_{l=0}^{N-1}W_N^{l(n-m)}=\begin{dcases}
            1, & \text{if } m-n\equiv 0 \mod N, \\
            0, & \otherwise.
        \end{dcases}
    \end{equation}
\end{proof}

如果不区分周期序列与有限长序列，则$\shah_N$就相当于周期序列中的单位冲激序列。用向量的语言来说，就是
\begin{equation}
    \shah_N\sim\boldsymbol{\delta}=\begin{bmatrix}
        1      \\
        0      \\
        \vdots \\
        0
    \end{bmatrix}.
\end{equation}

习惯上，将这种向量记法的各元素的序号依次取为$0,1,\dots,N-1$，而不是$1,2,\dots,N$，这样序号与离散傅里叶变换中的求和指标一致。仿照狄拉克分布的平移，离散选择函数的平移记为$\boldsymbol{\delta}_k[n]=\boldsymbol{\delta}[n-k]$，特别地，$\boldsymbol{\delta}_0=\boldsymbol{\delta}$。

类似狄拉克分布，容易验证单位脉冲序列具有取样性质和卷积性质：
\begin{proposition}
    设$x[n]$是周期为$N$的序列，则有
    \begin{align}
        x[n]\boldsymbol{\delta}_k[n]&=x[k]\boldsymbol{\delta}_k[n], \\
        x[n]*\boldsymbol{\delta}_k[n]&=x[n-k].
    \end{align}
\end{proposition}

\begin{definition}[序列的反转]
    将序列$x[n]$的反转序列$x[-n]$记为$x^-[n]$。将周期序列的反转序列写成向量形式，有：
    \begin{equation}
        \mathbf{x}^-=\begin{bmatrix}
            x[0]   \\
            x[N-1] \\
            x[N-2] \\
            \vdots \\
            x[2]   \\
            x[1]
        \end{bmatrix}.
    \end{equation}
\end{definition}

类似于连续函数的傅里叶变换，DFT的原序列与变换域序列都是周期的，此时可以研究DFT的对偶性。可以观察到卷积和定理与角频率形式的傅里叶变换的卷积定理类似，因此可以仿照角频率形式的傅里叶变换的对偶性写出DFT的对偶性的表述：
\begin{proposition}[离散傅里叶变换的对偶性]\label{prop:DFT_dual}
    设$x[n]$是周期为$N$的序列，其DFT为$X[k]$，则有
    \begin{align}
        \DFT(x^-)[k]&=X^-[k]=N\cdot\DFT^{-1}(x)[k], \\
        \DFT(\DFT(x))[n]&=N x^-[n].
    \end{align}
    如果$x[n]$是实序列，则有
    \begin{equation}
        \DFT(x^-)[k]=\overline{X[k]}=N\cdot\DFT^{-1}(x)[k].
    \end{equation}
\end{proposition}
\begin{proof}
    带入定义即可验证：
    \begin{align}
        \DFT(x^-)[k]&=\sum_{n=0}^{N-1}x^-[n]W_N^{kn} \nonumber \\
        &=\sum_{n=0}^{N-1}x[-n]W_N^{kn} \nonumber \\
        &=\sum_{m=0}^{N-1}x[m]W_N^{-km} \qquad (m=-n) \nonumber \\
        &=N\cdot\DFT^{-1}(x)[k],
    \end{align}
    以及
    \begin{align}
        \DFT(\DFT(x))[n]&=\sum_{k=0}^{N-1}X[k]W_N^{kn} \nonumber \\
        &=\sum_{k=0}^{N-1}\paren{\sum_{m=0}^{N-1}x[m]W_N^{km}}W_N^{kn} \nonumber \\
        &=\sum_{m=0}^{N-1}x[m]\paren{\sum_{k=0}^{N-1}W_N^{k(m+n)}} \nonumber \\
        &=N\sum_{m=0}^{N-1}x[m]\shah_N[m+n] \nonumber \\
        &=N x^-[n].
    \end{align}

    如果$x[n]$是实序列，则有
    \begin{align}
        \DFT(x^-)[k]&=\sum_{n=0}^{N-1}x^-[n]W_N^{kn} \nonumber \\
        &=\sum_{n=0}^{N-1}x[-n]W_N^{kn} \nonumber \\
        &=\overline{\sum_{n=0}^{N-1}x[-n]W_N^{-kn}} \nonumber \\
        &=\overline{X[k]}.
    \end{align}
\end{proof}

另外，DFT是一种\textbf{正交变换}。以$\mathcal{F}_N$表示周期为$N$的离散傅里叶变换对应的矩阵，则$\mathcal{F}_N$是正交矩阵，但不是单位正交矩阵：
\begin{equation}
    \mathcal{F}_N \mathcal{F}_N^H = N\cdot I_N,
\end{equation}
其中$H$为埃尔米特算子，表示取矩阵的共轭转置，$I_N$为$N$阶单位矩阵。正交变换保持向量的内积，即
\begin{equation}
    \forall \mathbf{x},\mathbf{y}\in\mathbb{C}^N,\quad \mathbf{x}\cdot\mathbf{y}=\frac{1}{N}\mathcal{F}_N \mathbf{x}\cdot\mathcal{F}_N \mathbf{y}.
\end{equation}
由此可得DFT的帕塞瓦尔定理：
\begin{theorem}[离散傅里叶变换的帕塞瓦尔定理]
    设$x[n],y[n]$是长为$N$的序列，其DFT分别为$X[k],Y[k]$，则有
    \begin{equation}
        \sum_{n=0}^{N-1}x[n]\overline{y[n]}=\frac{1}{N}\sum_{k=0}^{N-1}X[k]\overline{Y[k]}.
    \end{equation}
    特别地，当$x[n]=y[n]$时，有
    \begin{equation}
        \sum_{n=0}^{N-1}|x[n]|^2=\frac{1}{N}\sum_{k=0}^{N-1}|X[k]|^2.
    \end{equation}
\end{theorem}

\subsection{离散傅里叶变换与其他变换的关系}\label{sub:relation_of_dft}

\subsubsection{离散傅里叶变换可以看作是傅里叶级数展开的离散版本}

对于以$T$为周期的周期函数$f(t)\in L^2([0,T])$，记其傅里叶系数为$c_k$，如果$f(t)$连续，则有
\begin{equation}
    f(t)=\sum_{k=-\infty}^{\infty}c_k \e^{\im k\omega t}=\sum_{k=-\infty}^{\infty}c_k \e^{\im k\frac{2\pi}{T}t}.
\end{equation}

特别地，在一个周期内取$N$个等间隔点作为采样点，所得序列为
\begin{equation}
    x[n]=f\paren{\frac{nT}{N}},\quad n=0,1,\dots,N-1.
\end{equation}

对$f(t)$在采样点处的傅里叶级数进行分组求和，得到
\begin{align}
    x[n]=f\paren{\frac{nT}{N}}&=\sum_{k=-\infty}^{\infty}c_k \e^{\im k\frac{2\pi}{T}\frac{nT}{N}} \nonumber \\
    &=\sum_{k=-\infty}^{\infty}c_k W_N^{-kn} \nonumber \\
    &=\sum_{m=-\infty}^{\infty}\sum_{k=0}^{N-1}c_{k+mN} W_N^{-(k+mN)n} \nonumber \\
    &=\sum_{k=0}^{N-1}\paren{\sum_{m=-\infty}^{\infty}c_{k+mN}}W_N^{-kn}.
\end{align}

将上式与离散傅里叶反演公式进行比较，得到
\begin{equation}
    X[k]=N\sum_{m=-\infty}^{\infty}c_{k+mN}.
\end{equation}

这正是“\textbf{时域采样，频域周期化}”的观点在傅里叶级数上的体现。

\subsubsection{离散傅里叶变换可以看作离散时间傅里叶变换在时域上周期化、在频域上采样的结果}

设$x[n]$仅在$0\leq n\leq N-1$的范围内非零，它的DTFT为
\begin{equation}
    X(\e^{\im\Omega})=\sum_{n=0}^{N-1}x[n]\e^{-\im\Omega n}.
\end{equation}

将$x[n]$以$N$为周期进行周期化，得到周期为$N$的序列$\tilde{x}[n]$，则$\tilde{x}[n]$的DTFT为
\begin{align}
    \tilde{X}(\e^{\im\Omega})&=\sum_{n=-\infty}^{\infty}\tilde{x}[n]\e^{-\im\Omega n} \nonumber \\
    &=\sum_{m=-\infty}^{\infty}\sum_{n=0}^{N-1}x[n]\e^{-\im\Omega (n+mN)} \nonumber \\
    &=\paren{\sum_{m=-\infty}^{\infty}\e^{-\im\Omega mN}}\paren{\sum_{n=0}^{N-1}x[n]\e^{-\im\Omega n}}.
\end{align}
左侧括号内是狄拉克梳状分布的傅里叶级数形式（类似\cref{def:shah_fourier_series}），通过（角频率形式的）傅里叶变换可以推导出它的等价形式。由分布的傅里叶变换的连续性（\cref{thm:distribution_fourier_continuous}），有：
\begin{align}
    \mathcal{F}\sbr{\sum_{m=-\infty}^{\infty}\e^{-\im\Omega mN}}&=\sum_{m=-\infty}^{\infty}\mathcal{F}\sbr{\e^{-\im\Omega mN}} \nonumber \\
    &=2\pi\sum_{m=-\infty}^{\infty}\delta_{-mN} \nonumber \\
    &=2\pi\shah_{N}.
\end{align}
做傅里叶逆变换，得到：
\begin{equation}
    \sum_{m=-\infty}^{\infty}\e^{-\im\Omega mN}=2\pi\mathcal{F}^{-1}\shah_{N}=\frac{2\pi}{N}\shah_{2\pi/N},
\end{equation}
因此
\begin{equation}
    \tilde{X}(\e^{\im\Omega})=\frac{2\pi}{N}\shah_{2\pi/N}(\Omega)X(\e^{\im\Omega}).
\end{equation}

将$X(\e^{\im\Omega})$与$\shah_{2\pi/N}$相乘，相当于在$\Omega$轴上以$2\pi/N$为间隔进行采样，在点$\Omega_k=2\pi k/N$处，得到的样值为
\begin{equation}
    \sum_{n=0}^{N-1}x[n]\e^{-\im\Omega_k n}=\sum_{n=0}^{N-1}x[n]W_N^{-kn},
\end{equation}
这正是周期序列$\tilde{x}[n]$的DFT的值$X[k]$。

结合离散傅里叶反演公式可知：对于长为$N$的序列，只要给定了DTFT在$N$个等间隔点处的取值（这相当于对DTFT进行理想采样），就可以完全还原出原序列。

这正是“\textbf{时域周期化，频域采样}”的观点在DTFT上的体现。

如果由这个长为$N$的序列$x[n]$，定义一个长为$2N$的序列：
\begin{equation}
    y[n]=\begin{dcases}
        x[n], & 0 \leq n \leq N-1, \\
        0,    & N \leq n \leq 2N-1,
    \end{dcases}
\end{equation}
则$x[n]$与$y[n]$的DTFT是相同的。带入DFT的定义，容易验证：
\begin{equation}
    Y[2k]=X[k],\quad \forall k\in\{0,1,\dots,N-1\}.
\end{equation}

也就是说，在频域上，$Y[k]$相当于以两倍于$X[k]$的采样频率，对DTFT进行采样的结果，如下图所示：
\begin{figure}[H]
    \centering
    \includegraphics[width=0.8\textwidth]{discrete_sample}
    \caption{离散傅里叶变换与离散时间傅里叶变换的关系}
\end{figure}

上图中，$x[n]$的长度为$8$，$y[n]$的长度为$16$。随着采样点编号从小到大，所对应的数字角频率依次覆盖$0\to\pi\to2\pi$。因此，$|X(\e^{\im\Omega})|$在区间$[0,\pi]$上的曲线，对应于$|X[k]|$与$|Y[k]|$左半部分的灰色包络；而$|X(\e^{\im\Omega})|$在区间$[-\pi,0]$上的曲线，则对应于两者右半部分的灰色包络。

更一般地，通过延长$x[n]$并在新坐标点处取$0$，可以对$x[n]$的DTFT进行任意多采样点的等间隔采样，这就是\textbf{零填充}(Zero Padding)的思想。在需要区分DFT所针对的序列长度的情况下，我们使用\textbf{DFT点数}的概念，例如$X[k]$是$N$点DFT的结果，而$Y[k]$是$2N$点DFT的结果。

\subsubsection{离散傅里叶变换是连续傅里叶变换的逼近}

假设$f(t)$是在时域、频域上同时紧支的连续信号（\cref{thm:time_freq_limit}表明，这样的信号是不存在的，但近似满足这种关系的信号却很多），对$f(t)$进行理想采样时，将得到一个有限长序列$x[n]$\footnote{工程上，只有有限长的数字信号才能够存储和处理，但在信号与系统课程中，我们暂时不考虑量化带来的影响。}，在采样间隔趋向于$0$时，$x[n]$就能很好地还原$f(t)$的特征。下面说明，对$x[n]$的离散傅里叶变换$X[k]$在频域上也能很好地还原$f(t)$的特征。

简洁起见，假设$\supp(f)\subset[0,L]$，$\supp(\mathcal{F} f)\subset[-\omega_m,\omega_m]$，取$N$个等间隔点作为采样点，即
\begin{equation}
    x[n]=f\paren{\frac{nL}{N}},\quad n=0,1,\dots,N-1.
\end{equation}

采样间隔充分小时，采样频率高于奈奎斯特频率，即频域上周期化时不会产生交叠。因此，离散时间傅里叶变换$X(\e^{\im\Omega})$在$[-\pi,\pi]$范围内就是原函数的傅里叶变换$\mathcal{F} f(\omega)$在$[-\pi/T,\pi/T]$范围内的伸缩，即
\begin{equation}
    X(\e^{\im\Omega})=\mathcal{F} f\paren{\frac{\Omega}{T}},\quad \Omega\in[-\pi,\pi].
\end{equation}

因此，DTFT能够完全还原出原函数的频域特征。由于离散傅里叶变换是DTFT在频域上采样的结果，通过零填充的方式，总能使得离散傅里叶变换的频域采样点足够密集，从而很好地还原出原函数的频域特征。

\subsubsection{频谱泄露}

对于时域上无限长的带限函数，例如$f(t)=\cos(\omega_0 t)$，我们希望通过DFT对其进行频谱分析，此时应该观察到模拟角频域上$\omega=\pm\omega_0$处的两根谱线，在数字角频域作为模拟角频域的伸缩也应具有同样的形状，然而在使用DFT对其进行模拟时，效果可能很不好：

\begin{example}
    使用DFT对$f(t)=\cos(0.3\pi t)$进行频谱分析，则最高模拟角频率为$0.3\pi$，以$f_s=1\,\mathrm{s}$进行采样不会发生混叠，所得序列为
    \begin{equation}
        x[n]=f\paren{\frac{n}{f_s}}=\cos(0.3\pi n),
    \end{equation}
    $x[n]$是以$20$为周期的序列，我们提高DFT的点数$N$至$200$以便观察，所得结果如下图：
    \begin{figure}[H]
        \centering
        \includegraphics[width=0.6\textwidth]{spectral_leakage_1}
        \caption{200点离散傅里叶变换}
    \end{figure}
    可以看到两条清晰的线谱。

    下一小节中，我们将介绍DFT的快速算法FFT，它对$N$为$2$的幂次的情况表现最好也最简单。为此，我们尝试提高DFT的点数$N$至$256$，所得结果如下图：
    \begin{figure}[H]
        \centering
        \includegraphics[width=0.6\textwidth]{spectral_leakage}
        \caption{256点离散傅里叶变换}
    \end{figure}
    提高了DFT的点数，频谱反而出现了多余的谱线，这就是\textbf{频谱泄露}(spectral leakage)现象。尽管此时还能看到清晰的峰值，但当信号中有频率相近的成分时，频谱泄露就是不可接受的了。
\end{example}

细心的读者或许已经发现，当采用$256$点DFT时，$256$并不是$x[n]$的周期，这样进行DFT是不合理的，这就是发生频谱泄露的本质原因，工程上往往容易忽略这一点。

\subsection{快速傅里叶变换}\label{sub:FFT}

离散傅里叶变换处理的是有限长序列，能够直接在计算机上存储和计算，是我们使用计算机模拟其他变换的一个重要工具。离散傅里叶变换的直接计算需要$O(N^2)$的时间复杂度，但通过\textbf{快速傅里叶变换}(Fast Fourier Transform, FFT)算法，可以将时间复杂度降低到$O(N\log N)$，这使得离散傅里叶变换在实际应用中非常高效。本小节来介绍两种经典的快速傅里叶变换算法。

\subsubsection{基2时域抽选快速傅里叶变换(DIT-FFT)}

设$x[n]$是长为$N$的序列，$N=2^m$，将$x[n]$的离散傅里叶变换按$n$的奇偶性进行分组求和，得到
\begin{align}
    X[k]&=\sum_{n=0}^{N-1}x[n]W_N^{kn} \nonumber \\
    &=\sum_{n=0}^{N/2-1}x[2n]W_N^{k(2n)}+\sum_{n=0}^{N/2-1}x[2n+1]W_N^{k(2n+1)} \nonumber \\
    &=\sum_{n=0}^{N/2-1}x[2n]W_{N/2}^{kn}+W_N^{-k}\sum_{n=0}^{N/2-1}x[2n+1]W_{N/2}^{kn}.\label{eq:dit_fft_decomp}
\end{align}

$x[n]$的奇数项和偶数项都是长度为$N/2$的序列，他们的离散傅里叶变换分别为
\begin{align}
    X_1[k]&=\sum_{n=0}^{N/2-1}x[2n]W_{N/2}^{kn}, \\
    X_2[k]&=\sum_{n=0}^{N/2-1}x[2n+1]W_{N/2}^{kn}.
\end{align}
代回\refeqn{dit_fft_decomp}，并利用$W_N^{N/2}=-1$，得到
\begin{align}
    X[k]&=X_1[k]+W_N^{k}X_2[k], \\
    X[k+N/2]&=X_1[k]-W_N^{k}X_2[k].
\end{align}

因此，$X[k]$的计算可以通过计算$X_1[k]$和$X_2[k]$来实现，而$X_1[k]$和$X_2[k]$又是长度为$N/2$的序列的离散傅里叶变换。通过递归地将序列长度不断减半，最终可以将离散傅里叶变换的计算复杂度降低到$O(N\log N)$。

\subsubsection{基2频域抽选快速傅里叶变换(DIF-FFT)}

设$x[n]$是长为$N$的序列，$N=2^m$，将$x[n]$的离散傅里叶变换分为前后两组进行求和，得到
\begin{align}
    X[k]&=\sum_{n=0}^{N/2-1}x[n]W_N^{kn}+\sum_{n=N/2}^{N-1}x[n]W_N^{kn} \nonumber \\
    &=\sum_{n=0}^{N/2-1}x[n]W_N^{kn}+W_N^{kN/2}\sum_{n=0}^{N/2-1}x[n+N/2]W_N^{kn} \nonumber \\
    &=\sum_{n=0}^{N/2-1}x[n]W_N^{kn}+(-1)^k\sum_{n=0}^{N/2-1}x[n+N/2]W_N^{kn}.\label{eq:dif_fft_decomp}
\end{align}

$x[n]$的前半部分和后半部分都是长度为$N/2$的序列，他们的离散傅里叶变换分别为
\begin{align}
    X_1[k]&=\sum_{n=0}^{N/2-1}x[n]W_{N/2}^{kn}, \\
    X_2[k]&=\sum_{n=0}^{N/2-1}x[n+N/2]W_{N/2}^{kn}.
\end{align}
分别考虑$k$为偶数和奇数的情况，与\refeqn{dif_fft_decomp}对比，有
\begin{align}
    X[2k]&=\sum_{n=0}^{N/2-1}x[n]W_N^{2kn}+(-1)^{2k}\sum_{n=0}^{N/2-1}x[n+N/2]W_N^{2kn} \nonumber \\
    &=\sum_{n=0}^{N/2-1}x[n]W_{N/2}^{kn}+\sum_{n=0}^{N/2-1}x[n+N/2]W_{N/2}^{kn} \nonumber \\
    &=X_1[k]+X_2[k], \\
    X[2k+1]&=\sum_{n=0}^{N/2-1}x[n]W_N^{(2k+1)n}+(-1)^{2k+1}\sum_{n=0}^{N/2-1}x[n+N/2]W_N^{(2k+1)n} \nonumber \\
    &=W_{N}^n\sbr{\sum_{n=0}^{N/2-1}x[n]W_{N/2}^{kn}-\sum_{n=0}^{N/2-1}x[n+N/2]W_{N/2}^{kn}} \nonumber \\
    &=W_N^n(X_1[k]-X_2[k]).
\end{align}

与DIT-FFT相同，$X[k]$的计算可以通过计算$X_1[k]$和$X_2[k]$来实现，从而递归地将离散傅里叶变换的计算复杂度降低到$O(N\log N)$。

事实上，计算DIF-FFT的标准做法是对输入信号进行预处理。定义两个长为$N/2$的序列：
\begin{align}
    g_1[n]&=x[n]+x[n+N/2], \\
    g_2[n]&=(x[n]-x[n+N/2])W_N^n,
\end{align}
则根据DFT的线性性，有
\begin{align}
    X[2k]&=\DFT(g_1)[k]=\sum_{n=0}^{N/2-1}g_1[n]W_{N/2}^{kn}, \\
    X[2k+1]&=\DFT(g_2)[k]=\sum_{n=0}^{N/2-1}g_2[n]W_{N/2}^{kn}.
\end{align}

不难想到，离散傅里叶逆变换也能够通过类似的方式进行快速计算。

对于$N$不是$2$的幂次的情况，如果我们希望用DFT作为桥梁去近似计算其他变换的结果，那么可以通过零填充的方式将序列长度扩展到$2$的幂次，从而使用快速傅里叶变换算法进行计算。根据上一小节的内容，零填充能够使我们更好地逼近其他变换的结果。

如果$N$不是$2$的幂次，并且我们希望计算离散傅里叶变换本身，就需要使用其他快速算法。例如，当$N$可以分解为小质数的乘积时，可以进行混合基数快速傅里叶变换；当$N$是质数时，可以使用Bluestein算法或Chirp-Z变换算法；如果仅仅需要计算前几个频点的值，则可以使用Goertzel算法（见\cite{dsp}）。这些算法多数都已经被实现为高效的库函数，可以直接调用。